\documentclass[11pt]{article}
\usepackage[a4paper,margin=1in]{geometry}
\usepackage[T1]{fontenc}
\usepackage{amsmath}
\usepackage{booktabs}
\usepackage{enumitem}

\title{AI4MH: Credible Signal Detection for Mental Health Crisis Monitoring}
\author{Ujjwal Singh}
\date{}

\begin{document}
\maketitle

\section{Abstract}

Early detection of mental health crises using social media is limited by sparse data, noisy signals, and adversarial manipulation. Existing systems optimize classification accuracy, which fails under deployment conditions and produces unreliable escalation decisions.

This project proposes a credible signal detection framework that treats crisis monitoring as an uncertainty-aware inference problem. Instead of direct classification, the system estimates a Confidence-Weighted Crisis Index (chi) using sample stability, temporal consistency, and spatial corroboration.

A threat-filtering layer suppresses non-organic signals using lexical entropy, velocity anomalies, and diffusion patterns. The system emits confidence-calibrated escalation decisions with mandatory human review. The target outcome is decision reliability, not raw model accuracy.

\section{Problem Statement}

Current crisis-monitoring systems behave as classification pipelines. They fail in deployment for three structural reasons:

\begin{itemize}[leftmargin=*]
\item Data Sparsity: unstable predictions in low-activity regions.
\item Adversarial Noise: bot-driven or media-driven signal inflation.
\item Lack of Uncertainty: no confidence estimate attached to decisions.
\end{itemize}

\begin{quote}
How can a system estimate credible crisis signals under sparse, noisy, and adversarial conditions while producing calibrated confidence for human decision-making?
\end{quote}

\section{System Design}

\subsection{Overview}

The decision score combines inferred crisis likelihood, credibility, and suppression for adversarial activity.

\begin{equation*}
\text{Credible Signal} = \chi \times P(\text{crisis} \mid s) \times (1 - \text{threat})
\end{equation*}

\subsection{Architecture}

The pipeline is organized into four layers:

\begin{itemize}[leftmargin=*]
\item Data Ingestion
\item Signal Extraction
\item Credibility and Threat Modeling
\item Decision and Governance
\end{itemize}

Insert Figure 1 after this subsection. See Figure Production Prompts for the exact diagram specification.

\subsection{Credibility Model}

\begin{equation*}
\chi = \text{sample\_stability} \times \text{spatial\_support} \times \text{temporal\_consistency}
\end{equation*}

\begin{equation*}
\text{sample\_stability} = \min(1, \sqrt{N / k})
\end{equation*}

The credibility term penalizes low-volume regions, rewards agreement across neighboring regions, and rejects short-lived spikes that diverge from historical behavior.

Insert Figure 2 after this subsection. See Figure Production Prompts for the exact diagram specification.

\subsection{Threat Model}

\begin{itemize}[leftmargin=*]
\item Lexical entropy
\item URL density
\item Velocity anomaly
\end{itemize}

Threat scoring suppresses coordinated amplification, informational spikes, and other non-organic activity before any escalation is emitted.

\section{Implementation Plan}

\begin{itemize}[leftmargin=*]
\item Phase 1: Data ingestion
\item Phase 2: Signal extraction
\item Phase 3: Credibility modeling
\item Phase 4: Threat filtering
\item Phase 5: Decision engine
\item Phase 6: Evaluation
\item Phase 7: Governance
\end{itemize}

\section{Timeline}

\begin{tabular}{l l}
\toprule
Week & Task \\
\midrule
1-2 & Data pipeline \\
3 & Signal extraction \\
4 & Credibility modeling \\
5 & Threat filtering \\
6 & Decision engine \\
7 & Evaluation \\
8 & Governance \\
\bottomrule
\end{tabular}

\section{Deliverables}

\begin{itemize}[leftmargin=*]
\item Credible signal pipeline
\item Signal extraction module
\item Credibility model (chi)
\item Threat filter
\item Decision engine
\item Evaluation framework
\item Audit logging system
\end{itemize}

\section{Evaluation Plan}

\textbf{Primary Metric:} Precision @ 75\% Recall

\textbf{Secondary Metrics:} False Positive Rate, calibration error (ECE), and score stability.

Evaluation uses synthetic sparse-data, bot-amplification, media-spike, and organic-crisis scenarios to measure whether the system escalates only credible signals.

Insert Figure 3 in this section. See Figure Production Prompts for the exact scenario diagram specification.

\section{Risks and Mitigation}

\begin{itemize}[leftmargin=*]
\item Sparse Data: chi penalization and minimum-support gating.
\item Adversarial Attacks: entropy and velocity filters.
\item False Escalation: mandatory human-in-the-loop review.
\end{itemize}

\section{Why This Project}

The proposal targets decision reliability rather than classification accuracy. It aligns with AI4MH priorities around explainability, auditability, and governance while remaining deployable in public health monitoring settings.

\section{References}

\begin{enumerate}[leftmargin=*]
\item Atmakuru et al. (2025). AI-based suicide prevention and prediction: A systematic review. Information Fusion.
\item CLPsych shared tasks on suicide-risk detection in social media.
\item Moran's I for spatial autocorrelation and regional cluster analysis.
\item Bayesian inference and uncertainty-aware decision scoring methods.
\end{enumerate}

\section{Figure Production Prompts}

Limit the proposal to three figures. Do not add UI screenshots, generic machine-learning diagrams, or unrelated graphs.

\subsection*{Figure 1. System Architecture Diagram}
\textbf{Placement:} Place after System Design -> Architecture.

\begin{quote}
\ttfamily
Draw a clean pipeline diagram:\\\\Social Media Data ->\\Data Ingestion ->\\Signal Extraction ->\\Credibility Model (chi) ->\\Threat Filter ->\\Decision Engine ->\\Human Review\\\\Use rectangular blocks with arrows.\\Highlight chi and threat filter in bold.\\Minimal colors, academic style.
\end{quote}

\subsection*{Figure 2. Credible Signal Equation Diagram}
\textbf{Placement:} Place after System Design -> Credibility Model.

\begin{quote}
\ttfamily
Visualize equation:\\\\Credible Signal = chi x P(crisis \textbar{} s) x (1 - threat)\\\\Break into 3 components:\\1. chi (stability, spatial, temporal)\\2. P(crisis \textbar{} s)\\3. threat filter\\\\Use labeled blocks feeding into the final score.
\end{quote}

\subsection*{Figure 3. Evaluation Scenario Diagram}
\textbf{Placement:} Place in Evaluation Plan.

\begin{quote}
\ttfamily
Show 4 scenarios:\\\\1. Sparse data\\2. Bot attack\\3. Media spike\\4. Organic crisis\\\\Use a separate timeline plot for each scenario.\\Show expected system response: stable, suppressed, monitored, or detected.
\end{quote}

\end{document}
